% Vietnamese translation for OI Public Library
% Source-PDF: 比赛 CONTEST/2021“MINIEYE杯”中国大学生算法设计超级联赛/2021“MINIEYE杯”中国大学生算法设计超级联赛（3）/2021“MINIEYE杯”中国大学生算法设计超级联赛（3）-题目集.pdf
\documentclass[11pt]{article}
\usepackage{oipl-vn}

\title{Đề bài trận thứ ba HDU Multi-University 2021}
\author{}
\date{}

\begin{document}
\maketitle

\origtitle{Super League of Chinese College Students Algorithm Design 2021, Contest 3}
\sourcepath{contest/hdu-2021-minieye-03-problems.pdf}

\section{Problem A. Bookshop}

\paragraph{Tệp vào:} standard input
\paragraph{Tệp ra:} standard output
\paragraph{Giới hạn bộ nhớ:} 512 megabytes

Ngày mai DreamGrid sẽ đến hiệu sách. Trong hiệu sách có tổng cộng \(n\) cuốn
sách, được đánh số \(1,2,\ldots,n\), nối với nhau bởi \(n-1\) con đường hai
chiều tạo thành một cây. Vì DreamGrid rất giàu, cậu sẽ mua sách theo chiến
lược sau:
\begin{itemize}
  \item Chọn hai cuốn sách \(x\) và \(y\).
  \item Đi từ \(x\) đến \(y\) theo đường đi ngắn nhất trên cây, và lần lượt
  xem từng cuốn sách, bao gồm cả cuốn thứ \(x\) và cuốn thứ \(y\).
  \item Với mỗi cuốn sách đang được xem, nếu DreamGrid có đủ tiền (không ít
  hơn giá sách), cậu sẽ mua cuốn đó và số tiền còn lại giảm đi đúng bằng giá
  sách. Nếu số tiền của cậu nhỏ hơn giá của cuốn sách đang được xem, cậu sẽ
  bỏ qua cuốn sách đó.
\end{itemize}

BaoBao tò mò muốn biết DreamGrid giàu đến mức nào. Bạn được cho \(q\) truy
vấn; trong mỗi truy vấn có ba số nguyên \(x_i,y_i,z_i\). Giả sử DreamGrid
mang theo \(z_i\) đô la trước khi mua sách và đi từ \(x_i\) đến \(y_i\), hãy
cho BaoBao biết số tiền DreamGrid còn lại sau khi xem hết các cuốn sách đã
đi qua.

\subsection*{Dữ liệu vào}

Dòng đầu tiên chứa một số nguyên \(T\) (\(1\le T\le 500\)) là số bộ dữ liệu.
Với mỗi bộ dữ liệu:

Dòng đầu tiên chứa hai số nguyên \(n\) và \(q\)
(\(1\le n,q\le 100000\)), lần lượt là số cuốn sách trong hiệu sách và số
truy vấn.

Dòng thứ hai chứa \(n\) số nguyên \(p_1,p_2,\ldots,p_n\)
(\(1\le p_i\le 10^9\)), là giá của từng cuốn sách.

Mỗi dòng trong \(n-1\) dòng tiếp theo chứa hai số nguyên \(u_i\) và \(v_i\)
(\(1\le u_i,v_i\le n,\ u_i\ne v_i\)), biểu thị một con đường hai chiều giữa
\(u_i\) và \(v_i\). Bảo đảm rằng các con đường tạo thành một cây.

Trong \(q\) dòng tiếp theo, dòng thứ \(i\) chứa ba số nguyên \(x_i,y_i,z_i\)
(\(1\le x_i,y_i\le n,\ 1\le z_i\le 10^9\)), mô tả truy vấn thứ \(i\).

Bảo đảm rằng tổng tất cả các giá trị \(n\) không vượt quá \(800000\), và tổng
tất cả các giá trị \(q\) không vượt quá \(800000\).

\subsection*{Dữ liệu ra}

Với mỗi truy vấn, in một dòng chứa một số nguyên: số tiền DreamGrid còn lại
sau khi xem hết các cuốn sách đã đi qua.

\subsection*{Ví dụ}

\paragraph{standard input}
\begin{verbatim}
1
5 5
5 2 7 4 6
1 2
1 3
2 4
2 5
4 5 10
5 4 10
4 5 12
5 3 12
3 5 11
\end{verbatim}

\paragraph{standard output}
\begin{verbatim}
4
2
0
4
2
\end{verbatim}

\section{Problem B. Destinations}

\paragraph{Tệp vào:} standard input
\paragraph{Tệp ra:} standard output
\paragraph{Giới hạn bộ nhớ:} 512 megabytes

Có \(n\) thị trấn ở Byteland, được đánh số \(1,2,\ldots,n\), nối với nhau
bởi \(n-1\) con đường hai chiều tạo thành một cây. Có \(m\) du khách sắp tham
quan Byteland; du khách thứ \(i\) sẽ bắt đầu hành trình tại thị trấn
\(s_i\). Tuy nhiên, chưa ai quyết định nơi kết thúc hành trình. Cụ thể, du
khách thứ \(i\) đã lập \(3\) kế hoạch; kế hoạch thứ \(j\) kết thúc ở thị trấn
\(e_{i,j}\) và tốn \(c_{i,j}\) đô la. Khi một du khách đưa ra quyết định,
người đó sẽ đi theo đường đi ngắn nhất để thăm tất cả các thị trấn từ
\(s_i\) đến đích. Chú ý rằng \(e_{i,j}\) có thể trùng với \(s_i\); thị trấn
\(s_i\) và thị trấn đích cũng được du khách thứ \(i\) ghé thăm. Hai kế hoạch
có thể có cùng đích nhưng chi phí khác nhau, vì du khách có thể làm những
việc rất khác nhau tại cùng một thị trấn.

Các du khách sẽ chia sẻ ảnh về hành trình của mình ở Byteland. Không ai muốn
ghé thăm cùng một nơi với người khác. Nhiệm vụ của bạn là giúp mỗi du khách
chọn đích trong các kế hoạch của mình sao cho mỗi thị trấn được ghé thăm bởi
nhiều nhất một du khách, và tổng chi phí của tất cả du khách là nhỏ nhất;
hoặc xác định rằng điều này là không thể.

\subsection*{Dữ liệu vào}

Dòng đầu tiên chứa một số nguyên \(T\) (\(1\le T\le 500\)) là số bộ dữ liệu.
Với mỗi bộ dữ liệu:

Dòng đầu tiên chứa hai số nguyên \(n\) và \(m\)
(\(1\le n\le 200000,\ 1\le m\le 100000\)), lần lượt là số thị trấn và số du
khách.

Mỗi dòng trong \(n-1\) dòng tiếp theo chứa hai số nguyên \(u_i\) và \(v_i\)
(\(1\le u_i,v_i\le n,\ u_i\ne v_i\)), biểu thị một con đường hai chiều giữa
\(u_i\) và \(v_i\). Bảo đảm rằng các con đường tạo thành một cây.

Trong \(m\) dòng tiếp theo, dòng thứ \(i\) chứa bảy số nguyên
\(s_i,e_{i,1},c_{i,1},e_{i,2},c_{i,2},e_{i,3},c_{i,3}\)
(\(1\le s_i,e_{i,j}\le n,\ 1\le c_{i,j}\le 10^6\)), mô tả du khách thứ \(i\).

Bảo đảm rằng tổng tất cả các giá trị \(n\) không vượt quá \(1000000\), và
tổng tất cả các giá trị \(m\) không vượt quá \(300000\).

\subsection*{Dữ liệu ra}

Với mỗi bộ dữ liệu, in một dòng chứa một số nguyên là tổng chi phí nhỏ nhất.
Nếu không có nghiệm, in \texttt{-1}.

\subsection*{Ví dụ}

\paragraph{standard input}
\begin{verbatim}
2
7 2
1 2
1 3
2 4
2 5
3 6
3 7
2 1 1 3 100 4 200
3 2 1 4 2 7 50
4 2
1 2
2 3
3 4
1 2 1 3 1 4 1
2 1 1 3 1 4 1
\end{verbatim}

\paragraph{standard output}
\begin{verbatim}
51
-1
\end{verbatim}

\section{Problem C. Forgiving Matching}

\paragraph{Tệp vào:} standard input
\paragraph{Tệp ra:} standard output
\paragraph{Giới hạn bộ nhớ:} 512 megabytes

Little Q đang kiểm tra xem chuỗi \(A\) có khớp với chuỗi \(B\) hay không.
Hai chuỗi được xem là khớp nếu chúng có cùng độ dài và không có vị trí \(i\)
nào mà \(A_i\) khác \(B_i\).

Tuy nhiên, Little Q là người rộng lượng; cậu tha thứ cho mọi người từng làm
tổn thương mình. Hơn nữa, cậu còn tha thứ cho cả chuỗi! Cậu cho chuỗi \(k\)
cơ hội: nếu có không quá \(k\) vị trí \(i\) mà \(A_i\) khác \(B_i\), thì
Little Q vẫn xem hai chuỗi là khớp. Chú ý rằng cả hai chuỗi đều có thể chứa
ký tự đại diện \texttt{*}, ký tự này có thể khớp đúng một ký tự bất kỳ; trong
trường hợp đó, cặp ký tự này không tiêu tốn cơ hội tha thứ.

Ký hiệu \(\operatorname{occ}(S,T)\) là số chuỗi con trong chuỗi \(S\) khớp
với \(T\); hai chuỗi con được xem là khác nhau nếu chúng bắt đầu ở hai vị trí
khác nhau. Bạn được cho hai chuỗi \(S\) và \(T\). Hãy viết chương trình tính
giá trị \(\operatorname{occ}(S,T)\) với
\(k=0,1,2,\ldots,|T|\).

\subsection*{Dữ liệu vào}

Dòng đầu tiên chứa một số nguyên \(K\) (\(1\le K\le 100\)) là số bộ dữ liệu.
Với mỗi bộ dữ liệu:

Dòng đầu tiên chứa hai số nguyên \(n\) và \(m\)
(\(1\le m\le n\le 200000\)), lần lượt là độ dài của \(S\) và độ dài của \(T\).

Dòng thứ hai chứa chuỗi \(S\) độ dài \(n\).

Dòng thứ ba chứa chuỗi \(T\) độ dài \(m\).

Bảo đảm rằng tổng tất cả các giá trị \(n\) không vượt quá \(1000000\). Cả
\(S\) và \(T\) chỉ có thể chứa các ký tự trong
\(\{\texttt{0},\texttt{1},\texttt{2},\texttt{3},\texttt{4},\texttt{5},
\texttt{6},\texttt{7},\texttt{8},\texttt{9},\texttt{*}\}\).

\subsection*{Dữ liệu ra}

Với mỗi bộ dữ liệu, in \(m+1\) dòng. Dòng thứ \(i\)
(\(1\le i\le m+1\)) chứa một số nguyên, là giá trị
\(\operatorname{occ}(S,T)\) khi \(k=i-1\).

\subsection*{Ví dụ}

\paragraph{standard input}
\begin{verbatim}
1
5 3
012*4
1*3
\end{verbatim}

\paragraph{standard output}
\begin{verbatim}
1
1
3
3
\end{verbatim}

\section{Problem D. Game on Plane}

\paragraph{Tệp vào:} standard input
\paragraph{Tệp ra:} standard output
\paragraph{Giới hạn bộ nhớ:} 512 megabytes

Alice và Bob đang chơi một trò chơi. Trong trò chơi này có \(n\) đường thẳng
trên mặt phẳng hai chiều. Trước hết Alice sẽ chọn đúng \(k\) đường thẳng
\(l_1,l_2,\ldots,l_k\) trong số \(n\) đường thẳng, sau đó Bob sẽ vẽ một
đường thẳng \(L\). Hình phạt của Bob được định nghĩa là số đường thẳng trong
\(\{l_1,l_2,\ldots,l_k\}\) có ít nhất một điểm chung với \(L\). Chú ý rằng
hai đường thẳng trùng nhau cũng có các điểm chung.

Alice muốn tối đa hóa hình phạt của Bob, còn Bob muốn tối thiểu hóa nó. Bạn
được cho \(n\) đường thẳng này; hãy viết chương trình dự đoán hình phạt của
Bob với \(k=1,2,3,\ldots,n\), nếu cả hai người đều chơi tối ưu.

\subsection*{Dữ liệu vào}

Dòng đầu tiên chứa một số nguyên \(T\) (\(1\le T\le 500\)) là số bộ dữ liệu.
Với mỗi bộ dữ liệu:

Dòng đầu tiên chứa một số nguyên \(n\) (\(1\le n\le 100000\)), là số đường
thẳng.

Mỗi dòng trong \(n\) dòng tiếp theo chứa bốn số nguyên
\(x_{a_i},y_{a_i},x_{b_i},y_{b_i}\)
(\(0\le x_{a_i},y_{a_i},x_{b_i},y_{b_i}\le 10^9\)), biểu thị một đường
thẳng đi qua cả \((x_{a_i},y_{a_i})\) và \((x_{b_i},y_{b_i})\).
\((x_{a_i},y_{a_i})\) sẽ không bao giờ trùng với \((x_{b_i},y_{b_i})\).

Bảo đảm rằng tổng tất cả các giá trị \(n\) không vượt quá \(1000000\).

\subsection*{Dữ liệu ra}

Với mỗi bộ dữ liệu, in \(n\) dòng. Dòng thứ \(i\) (\(1\le i\le n\)) chứa một
số nguyên, là hình phạt của Bob khi \(k=i\).

\subsection*{Ví dụ}

\paragraph{standard input}
\begin{verbatim}
2
2
1 1 2 2
0 0 2 3
3
1 1 2 2
1 1 2 2
3 2 5 4
\end{verbatim}

\paragraph{standard output}
\begin{verbatim}
0
1
0
0
0
\end{verbatim}

\section{Problem E. Kart Race}

\paragraph{Tệp vào:} standard input
\paragraph{Tệp ra:} standard output
\paragraph{Giới hạn bộ nhớ:} 512 megabytes

Cuộc đua xe kart thường niên sắp được tổ chức ở Byteland. Bản đồ đường đua
gồm \(n\) giao lộ khác nhau và \(m\) đường một chiều. Các giao lộ được đánh
số từ \(1\) đến \(n\); giao lộ thứ \(i\) nằm tại \((x_i,y_i)\). Mạng lưới
đường không có chu trình; người ta chỉ có thể lái đến các giao lộ có tọa độ
\(x\) lớn hơn nghiêm ngặt. Các con đường chỉ có thể cắt nhau tại các giao lộ.

Cuộc đua bắt đầu ở giao lộ thứ \(1\) và kết thúc ở giao lộ thứ \(n\). Các tay
đua có thể tự chọn tuyến đường, nhưng chỉ được đi theo các con đường được
đánh dấu trên bản đồ. Bảo đảm rằng từ \(1\) có thể đi đến mọi nơi, và từ mọi
nơi có thể đi đến \(n\), nên mọi tuyến đường đều hợp lệ.

Cuộc đua xe kart thu hút rất nhiều nhà tài trợ. Mỗi giao lộ có một vị trí để
đặt băng rôn; nếu bạn chọn đặt băng rôn ở giao lộ thứ \(i\), công ty tổ chức
đua sẽ thu được \(w_i\) lợi nhuận. Bạn là người trung gian trong công ty tổ
chức đua; nhiệm vụ của bạn là chọn một số giao lộ để đặt băng rôn sao cho
tổng lợi nhuận là lớn nhất. Bạn biết rằng không tay đua nào muốn nhìn thấy
nhiều hơn một băng rôn, nên với mọi tuyến đường có thể từ \(1\) đến \(n\),
bạn phải bảo đảm rằng nhiều nhất một giao lộ được chọn.

\subsection*{Dữ liệu vào}

Dòng đầu tiên chứa một số nguyên \(T\) (\(1\le T\le 1000\)) là số bộ dữ liệu.
Với mỗi bộ dữ liệu:

Dòng đầu tiên chứa hai số nguyên \(n\) và \(m\)
(\(1\le n\le 100000,\ 1\le m\le 2n\)), lần lượt là số giao lộ và số đường một
chiều.

Trong \(n\) dòng tiếp theo, dòng thứ \(i\) chứa ba số nguyên \(x_i,y_i,w_i\)
(\(0\le x_i,y_i\le 10^9,\ 1\le w_i\le 10^9\)), mô tả giao lộ thứ \(i\). Bảo
đảm rằng không có hai giao lộ nào có cùng tọa độ.

Mỗi dòng trong \(m\) dòng tiếp theo chứa hai số nguyên \(u_i\) và \(v_i\)
(\(1\le u_i,v_i\le n,\ x_{u_i}<x_{v_i}\)), biểu thị một con đường một chiều
từ \(u_i\) đến \(v_i\). Bảo đảm rằng mỗi cặp \(u_i\) và \(v_i\) được mô tả
nhiều nhất một lần.

Bảo đảm rằng tổng tất cả các giá trị \(n\) không vượt quá \(1500000\).

\subsection*{Dữ liệu ra}

Với mỗi bộ dữ liệu, in một số nguyên ở dòng đầu tiên, là tổng lợi nhuận lớn
nhất. Sau đó in một dãy số nguyên ở dòng thứ hai, là chỉ số các giao lộ bạn
chọn để đặt băng rôn. Nếu có nhiều nghiệm tối ưu, hãy in nghiệm nhỏ nhất theo
thứ tự từ điển.

\subsection*{Ví dụ}

\paragraph{standard input}
\begin{verbatim}
2
6 6
0 1 1
2 2 1
1 0 1
1 2 1
2 0 1
3 1 1
1 4
3 5
2 6
5 6
1 3
4 2
2 1
0 0 8
1 1 9
1 2
\end{verbatim}

\paragraph{standard output}
\begin{verbatim}
2
2 3
9
2
\end{verbatim}

\section{Problem F. New Equipments II}

\paragraph{Tệp vào:} standard input
\paragraph{Tệp ra:} standard output
\paragraph{Giới hạn bộ nhớ:} 512 megabytes

Nhà máy của Little Q gần đây mua \(n\) thiết bị mới, được đánh số
\(1,2,\ldots,n\).

Trong nhà máy có \(n\) công nhân, được đánh số \(1,2,\ldots,n\). Mỗi công
nhân có thể được phân công cho không quá một thiết bị, và không thiết bị nào
có thể được phân công cho nhiều công nhân. Nếu Little Q phân công công nhân
thứ \(i\) cho thiết bị thứ \(j\), họ sẽ đem lại \(a_i+b_j\) lợi nhuận. Tuy
nhiên, các công nhân này chưa có nhiều kinh nghiệm, nên có thêm \(m\) ràng
buộc. Trong mỗi ràng buộc, bạn được cho hai số nguyên \(u_i\) và \(v_i\),
biểu thị công nhân thứ \(u_i\) không thể được phân công cho thiết bị thứ
\(v_i\).

Bây giờ, với mỗi \(k\) (\(1\le k\le n\)), hãy tìm \(k\) cặp công nhân và
thiết bị rồi phân công công nhân cho các thiết bị đó sao cho tổng lợi nhuận
của \(k\) cặp là lớn nhất, hoặc xác định rằng điều này là không thể.

\subsection*{Dữ liệu vào}

Dòng đầu tiên chứa một số nguyên \(T\) (\(1\le T\le 200\)) là số bộ dữ liệu.
Với mỗi bộ dữ liệu:

Dòng đầu tiên chứa hai số nguyên \(n\) và \(m\)
(\(1\le n\le 4000,\ 1\le m\le 10000\)), lần lượt là số công nhân/thiết bị mới
và số ràng buộc.

Dòng thứ hai chứa \(n\) số nguyên \(a_1,a_2,\ldots,a_n\)
(\(1\le a_i\le 10^9\)).

Dòng thứ ba chứa \(n\) số nguyên \(b_1,b_2,\ldots,b_n\)
(\(1\le b_i\le 10^9\)).

Mỗi dòng trong \(m\) dòng tiếp theo chứa hai số nguyên \(u_i\) và \(v_i\)
(\(1\le u_i,v_i\le n\)), biểu thị công nhân thứ \(u_i\) không thể được phân
công cho thiết bị thứ \(v_i\). Mỗi cặp \(u_i\) và \(v_i\) được mô tả nhiều
nhất một lần.

Bảo đảm rằng có nhiều nhất \(10\) bộ dữ liệu có \(n>100\).

\subsection*{Dữ liệu ra}

Với mỗi bộ dữ liệu, in \(n\) dòng. Dòng thứ \(k\) (\(1\le k\le n\)) chứa một
số nguyên, là tổng lợi nhuận lớn nhất có thể đạt được cho \(k\) cặp công nhân
và thiết bị. Nếu không thể tìm được \(k\) cặp như vậy, in \texttt{-1} ở dòng
này.

\subsection*{Ví dụ}

\paragraph{standard input}
\begin{verbatim}
2
4 4
10 20 30 40
5 10 7 8
4 2
3 2
1 4
1 2
2 2
1 1
1 1
1 1
1 2
\end{verbatim}

\paragraph{standard output}
\begin{verbatim}
48
85
115
130
2
-1
\end{verbatim}

\section{Problem G. Photoshop Layers}

\paragraph{Tệp vào:} standard input
\paragraph{Tệp ra:} standard output
\paragraph{Giới hạn bộ nhớ:} 512 megabytes

Các điểm ảnh trong một bức ảnh số có thể được biểu diễn bằng ba số nguyên
\((R,G,B)\) trong khoảng từ \(0\) đến \(255\), cho biết cường độ của các màu
đỏ, xanh lá và xanh dương. Màu của một điểm ảnh có thể được biểu diễn bằng
một chuỗi chữ hoa thập lục phân gồm sáu chữ số. Ví dụ,
\((R=100,G=255,B=50)\) có thể được biểu diễn là \texttt{64FF32}.

Trong vùng làm việc Photoshop có \(n\) lớp, được đánh số \(1,2,\ldots,n\) từ
dưới lên trên. Màn hình sẽ hiển thị các lớp này theo thứ tự từ dưới lên trên.
Trong bài này, bạn chỉ cần xử lý trường hợp mọi điểm ảnh trong một lớp đều có
cùng màu. Màu của lớp thứ \(i\) là \(c_i=(R_i,G_i,B_i)\), và chế độ hòa trộn
của lớp thứ \(i\) là \(m_i\) (\(m_i\in\{1,2\}\)):
\begin{itemize}
  \item Nếu \(m_i=1\), chế độ hòa trộn của lớp này là ``Normal''. Giả sử màu
  đang hiển thị trước đó trên màn hình là \((R_p,G_p,B_p)\), màu mới sẽ là
  \((R_i,G_i,B_i)\).
  \item Nếu \(m_i=2\), chế độ hòa trộn của lớp này là ``Linear Dodge''. Giả
  sử màu đang hiển thị trước đó trên màn hình là \((R_p,G_p,B_p)\), màu mới
  sẽ là
  \[
    (\min(R_p+R_i,255),\min(G_p+G_i,255),\min(B_p+B_i,255)).
  \]
\end{itemize}

Bạn được cho \(q\) truy vấn. Trong truy vấn thứ \(i\), bạn được cho hai số
nguyên \(l_i\) và \(r_i\) (\(1\le l_i\le r_i\le n\)). Hãy viết chương trình
tính màu cuối cùng hiển thị trên màn hình nếu chỉ giữ lại tất cả các lớp có
chỉ số trong \([l_i,r_i]\) mà không thay đổi thứ tự của chúng. Chú ý rằng màu
nền là \((R=0,G=0,B=0)\).

\subsection*{Dữ liệu vào}

Dòng đầu tiên chứa một số nguyên \(T\) (\(1\le T\le 10\)) là số bộ dữ liệu.
Với mỗi bộ dữ liệu:

Dòng đầu tiên chứa hai số nguyên \(n\) và \(q\)
(\(1\le n,q\le 100000\)), lần lượt là số lớp và số truy vấn.

Trong \(n\) dòng tiếp theo, dòng thứ \(i\) chứa một số nguyên \(m_i\) và một
chuỗi chữ hoa thập lục phân gồm sáu chữ số \(c_i\), mô tả lớp thứ \(i\).

Trong \(q\) dòng tiếp theo, dòng thứ \(i\) chứa hai số nguyên \(l_i\) và
\(r_i\) (\(1\le l_i\le r_i\le n\)), mô tả truy vấn thứ \(i\).

\subsection*{Dữ liệu ra}

Với mỗi truy vấn, in một dòng chứa một chuỗi chữ hoa thập lục phân gồm sáu
chữ số, là màu cuối cùng được hiển thị.

\subsection*{Ví dụ}

\paragraph{standard input}
\begin{verbatim}
1
5 5
1 64C832
2 000100
2 010001
1 323C21
2 32C8C8
1 2
1 3
2 3
2 4
2 5
\end{verbatim}

\paragraph{standard output}
\begin{verbatim}
64C932
65C933
010101
323C21
64FFE9
\end{verbatim}

\section{Problem H. Restore Atlantis II}

\paragraph{Tệp vào:} standard input
\paragraph{Tệp ra:} standard output
\paragraph{Giới hạn bộ nhớ:} 512 megabytes

Có \(n\) bản đồ Hy Lạp cổ đại mô tả quần đảo huyền thoại Atlantis. Các bản đồ
được đánh số \(1,2,\ldots,n\). Bản đồ thứ \(i\) cho biết vùng hình chữ nhật
\(R_i\) là một phần của Atlantis. Các cạnh của mọi hình chữ nhật đều song
song với các trục. Có thể có nhiều đảo, và các hình chữ nhật có thể chồng
lên nhau.

Không may là một số bản đồ không đáng tin nên sẽ không được xét. Bạn được cho
\(q\) truy vấn. Trong truy vấn thứ \(i\), bạn được cho hai số nguyên \(s_i\)
và \(t_i\) (\(1\le s_i\le t_i\le n\)). Hãy viết chương trình tính tổng diện
tích của Atlantis khi chỉ các bản đồ có nhãn \(k\) (\(s_i\le k\le t_i\)) là
đáng tin.

\subsection*{Dữ liệu vào}

Dòng đầu tiên chứa một số nguyên \(T\) (\(1\le T\le 3\)) là số bộ dữ liệu.
Với mỗi bộ dữ liệu:

Dòng đầu tiên chứa hai số nguyên \(n\) và \(q\)
(\(1\le n,q\le 100000\)), lần lượt là số bản đồ và số truy vấn.

Trong \(n\) dòng tiếp theo, dòng thứ \(i\) chứa bốn số nguyên
\(x_{a_i},y_{a_i},x_{b_i},y_{b_i}\)
(\(0\le x_{a_i}<x_{b_i}\le 10^9,\ 0\le y_{a_i}<y_{b_i}\le 10^9\)), mô tả
bản đồ thứ \(i\), \(R_i\). \((x_{a_i},y_{a_i})\) là góc tây nam của \(R_i\),
và \((x_{b_i},y_{b_i})\) là góc đông bắc của \(R_i\).

Trong \(q\) dòng tiếp theo, dòng thứ \(i\) chứa hai số nguyên \(s_i\) và
\(t_i\) (\(1\le s_i\le t_i\le n\)), mô tả truy vấn thứ \(i\).

Bảo đảm rằng mọi giá trị \(x_{a_i},y_{a_i},x_{b_i},y_{b_i},s_i,t_i\) được
chọn ngẫu nhiên đều từ các số nguyên trong khoảng tương ứng của chúng. Điều
kiện ngẫu nhiên không áp dụng cho bộ dữ liệu mẫu, nhưng lời giải của bạn vẫn
phải vượt qua mẫu.

\subsection*{Dữ liệu ra}

Với mỗi truy vấn, in một dòng chứa một số nguyên, là tổng diện tích thu được
từ thông tin của tất cả các bản đồ đáng tin trong truy vấn này.

\subsection*{Ví dụ}

\paragraph{standard input}
\begin{verbatim}
1
3 6
10 10 20 20
12 12 22 22
15 15 25 25
1 1
2 2
3 3
1 2
2 3
1 3
\end{verbatim}

\paragraph{standard output}
\begin{verbatim}
100
100
100
136
151
187
\end{verbatim}

\section{Problem I. Rise in Price}

\paragraph{Tệp vào:} standard input
\paragraph{Tệp ra:} standard output
\paragraph{Giới hạn bộ nhớ:} 512 megabytes

Có \(n\times n\) ô trên một lưới; ô góc trên bên trái là \((1,1)\), còn ô
góc dưới bên phải là \((n,n)\). Bạn bắt đầu ở \((1,1)\) và di chuyển đến
\((n,n)\). Tại bất kỳ ô \((i,j)\), bạn có thể đi đến \((i+1,j)\) hoặc
\((i,j+1)\), miễn là không đi ra ngoài lưới. Rõ ràng bạn sẽ đi đúng
\(2n-2\) bước.

Khi ở ô \((i,j)\), bao gồm cả điểm bắt đầu \((1,1)\) và đích \((n,n)\), bạn
có thể lấy tất cả \(a_{i,j}\) viên kim cương tại ô này, và có cơ hội tăng
giá của mỗi viên kim cương thêm \(b_{i,j}\) đô la. Cuối cùng bạn sẽ bán tất
cả kim cương mình có với mức giá cuối cùng, và mục tiêu của bạn là chọn đường
đi tối ưu để tối đa hóa lợi nhuận. Chú ý rằng ban đầu giá của mỗi viên kim
cương bằng không, và bạn không có gì để bán.

\subsection*{Dữ liệu vào}

Dòng đầu tiên chứa một số nguyên \(T\) (\(1\le T\le 10\)) là số bộ dữ liệu.
Với mỗi bộ dữ liệu:

Dòng đầu tiên chứa một số nguyên \(n\) (\(1\le n\le 100\)), là kích thước
của lưới.

Mỗi dòng trong \(n\) dòng tiếp theo chứa \(n\) số nguyên; dòng thứ \(i\) chứa
\(a_{i,1},a_{i,2},\ldots,a_{i,n}\) (\(1\le a_{i,j}\le 10^6\)), là số viên
kim cương trong từng ô.

Mỗi dòng trong \(n\) dòng tiếp theo chứa \(n\) số nguyên; dòng thứ \(i\) chứa
\(b_{i,1},b_{i,2},\ldots,b_{i,n}\) (\(1\le b_{i,j}\le 10^6\)), là số tiền
bạn có thể tăng giá tại từng ô.

Bảo đảm rằng mọi giá trị \(a_{i,j}\) và \(b_{i,j}\) được chọn ngẫu nhiên đều
từ các số nguyên trong \([1,10^6]\). Điều kiện ngẫu nhiên không áp dụng cho
bộ dữ liệu mẫu, nhưng lời giải của bạn vẫn phải vượt qua mẫu.

\subsection*{Dữ liệu ra}

Với mỗi bộ dữ liệu, in một dòng chứa một số nguyên: số đô la lớn nhất bạn có
thể kiếm được bằng cách bán kim cương.

\subsection*{Ví dụ}

\paragraph{standard input}
\begin{verbatim}
1
4
2 3 1 5
6 3 2 4
3 5 1 4
5 2 4 1
3 2 5 1
2 4 3 5
1 2 3 4
4 3 5 3
\end{verbatim}

\paragraph{standard output}
\begin{verbatim}
528
\end{verbatim}

\section{Problem J. Road Discount}

\paragraph{Tệp vào:} standard input
\paragraph{Tệp ra:} standard output
\paragraph{Giới hạn bộ nhớ:} 512 megabytes

Có \(n\) thành phố ở Byteland, được đánh số từ \(1\) đến \(n\). Cơ quan xây
dựng giao thông của Byteland đang lập kế hoạch xây \(n-1\) con đường hai
chiều giữa các thành phố này sao cho mọi cặp thành phố khác nhau đều được
kết nối trực tiếp hoặc gián tiếp bởi các con đường đó.

Công ty kỹ thuật đã đề xuất \(m\) con đường ứng viên có thể xây dựng. Con
đường ứng viên thứ \(i\) tốn \(c_i\) đô la; nếu cuối cùng được xây, nó sẽ nối
trực tiếp thành phố thứ \(u_i\) và thành phố thứ \(v_i\). May mắn là mỗi con
đường có giá giảm riêng, giá giảm của con đường thứ \(i\) là \(d_i\).

Cơ quan xây dựng giao thông của Byteland có thể mua nhiều nhất \(k\) con
đường với giá giảm. Hãy viết chương trình giúp cơ quan này tìm phương án rẻ
nhất với \(k=0,1,2,\ldots,n-1\).

\subsection*{Dữ liệu vào}

Dòng đầu tiên chứa một số nguyên \(T\) (\(1\le T\le 10\)) là số bộ dữ liệu.
Với mỗi bộ dữ liệu:

Dòng đầu tiên chứa hai số nguyên \(n\) và \(m\)
(\(2\le n\le 1000,\ n-1\le m\le 200000\)), lần lượt là số thành phố và số con
đường ứng viên.

Mỗi dòng trong \(m\) dòng tiếp theo chứa bốn số nguyên \(u_i,v_i,c_i,d_i\)
(\(1\le u_i,v_i\le n,\ u_i\ne v_i,\ 1\le d_i\le c_i\le 1000\)), mô tả một
con đường ứng viên.

\subsection*{Dữ liệu ra}

Với mỗi bộ dữ liệu, in \(n\) dòng. Dòng thứ \(i\) (\(1\le i\le n\)) chứa một
số nguyên, là tổng chi phí rẻ nhất để xây \(n-1\) con đường khi \(k=i-1\).
Bảo đảm rằng đáp án luôn tồn tại.

\subsection*{Ví dụ}

\paragraph{standard input}
\begin{verbatim}
1
5 6
1 2 1 1
2 3 2 1
2 4 3 2
2 5 4 3
1 3 5 3
4 5 6 1
\end{verbatim}

\paragraph{standard output}
\begin{verbatim}
10
7
6
5
5
\end{verbatim}

\section{Problem K. Segment Tree with Pruning}

\paragraph{Tệp vào:} standard input
\paragraph{Tệp ra:} standard output
\paragraph{Giới hạn bộ nhớ:} 512 megabytes

Chenjb hiện đang vật lộn với cấu trúc dữ liệu. Cậu đang cố giải một bài toán
bằng cây đoạn. Chenjb là tân binh trong lập trình thi đấu, và cậu viết đoạn
mã C/C++ sau rồi chạy \texttt{Node* root = build(1, n)} để xây một cây đoạn
chuẩn trên đoạn \([1,n]\):

\begin{verbatim}
1 Node * build ( long long l , long long r ) {
2      Node * x = new ( Node );
3      if ( l == r ) return x ;
4      long long mid = ( l + r ) / 2;
5      x -> lchild = build (l , mid );
6      x -> rchild = build ( mid + 1 , r );
7      return x ;
8 }
\end{verbatim}

Chenjb nộp mã, nhưng không may bị MLE (vượt quá giới hạn bộ nhớ). Chẳng bao
lâu sau Chenjb nhận ra chương trình của mình sẽ cấp phát mới rất nhiều nút,
và cậu quyết định giảm số nút bằng cách cắt tỉa:

\begin{verbatim}
1 Node * build ( long long l , long long r ) {
2      Node * x = new ( Node );
3      if ( r - l + 1 <= k ) return x ;
4      long long mid = ( l + r ) / 2;
5      x -> lchild = build (l , mid );
6      x -> rchild = build ( mid + 1 , r );
7      return x ;
8 }
\end{verbatim}

Bạn biết rằng Chenjb là tân binh, nên cậu sẽ thử các giá trị \(k\) khác nhau
để tìm giá trị tối ưu. Bạn được cho các giá trị \(n\) và \(k\); hãy cho cậu
biết số nút được sinh ra bởi chương trình mới của cậu.

\subsection*{Dữ liệu vào}

Dòng đầu tiên chứa một số nguyên \(T\) (\(1\le T\le 10000\)) là số bộ dữ
liệu. Với mỗi bộ dữ liệu:

Dòng duy nhất chứa hai số nguyên \(n\) và \(k\)
(\(1\le k\le n\le 10^{18}\)), mô tả một truy vấn.

\subsection*{Dữ liệu ra}

Với mỗi truy vấn, in một dòng chứa một số nguyên, là số nút của cây đoạn.

\subsection*{Ví dụ}

\paragraph{standard input}
\begin{verbatim}
3
100000 1
100000 50
1000000000000000000 1
\end{verbatim}

\paragraph{standard output}
\begin{verbatim}
199999
4095
1999999999999999999
\end{verbatim}

\section{Problem L. Tree Planting}

\paragraph{Tệp vào:} standard input
\paragraph{Tệp ra:} standard output
\paragraph{Giới hạn bộ nhớ:} 512 megabytes

Có \(n\) tòa nhà bên đường Bytestreet, đứng liên tiếp cạnh nhau, được đánh số
\(1,2,\ldots,n\) từ đông sang tây. Đoàn khảo sát quy hoạch đô thị sắp trồng
một số cây trước các tòa nhà này. Một kế hoạch trồng cây được gọi là đẹp khi
và chỉ khi nó thỏa mãn tất cả các điều kiện sau:
\begin{itemize}
  \item Có ít nhất một cây được trồng.
  \item Một cây chỉ có thể được trồng trước một tòa nhà, và hai cây không thể
  được trồng trước cùng một tòa nhà. Nói cách khác, nếu \(c_i\) là số cây
  được trồng trước tòa nhà thứ \(i\), thì \(c_i\in\{0,1\}\) với mọi
  \(i=1,2,\ldots,n\).
  \item Với mọi giá trị \(i\) có thể (\(1\le i<n\)), \(c_i\) và \(c_{i+1}\)
  không được đồng thời bằng \(1\).
  \item Với mọi giá trị \(i\) có thể (\(1\le i\le n-k\)), \(c_i\) và
  \(c_{i+k}\) không được đồng thời bằng \(1\).
\end{itemize}

Độ đẹp của một kế hoạch được định nghĩa là
\[
  \prod_{1\le i\le n,\ c_i=1} w_i.
\]

Bạn cần tìm tổng độ đẹp trên tất cả các kế hoạch đẹp. Hai kế hoạch được xem
là khác nhau nếu tồn tại một \(i\) (\(1\le i\le n\)) sao cho chúng khác nhau
ở \(c_i\).

\subsection*{Dữ liệu vào}

Dòng đầu tiên chứa một số nguyên \(T\) (\(1\le T\le 200\)) là số bộ dữ liệu.
Với mỗi bộ dữ liệu:

Dòng đầu tiên chứa hai số nguyên \(n\) và \(k\)
(\(2\le k<n\le 300\)), lần lượt là số tòa nhà và giá trị \(k\).

Dòng thứ hai chứa \(n\) số nguyên \(w_1,w_2,\ldots,w_n\)
(\(1\le w_i\le 10^9\)).

Bảo đảm rằng có nhiều nhất \(10\) bộ dữ liệu có \(n>100\).

\subsection*{Dữ liệu ra}

Với mỗi bộ dữ liệu, in một dòng chứa một số nguyên, là tổng độ đẹp trên tất
cả các kế hoạch đẹp. Chú ý rằng đáp án có thể cực kỳ lớn, vì vậy hãy in nó
theo modulo \(10^9+7\).

\subsection*{Ví dụ}

\paragraph{standard input}
\begin{verbatim}
2
5 3
1 1 1 1 1
5 4
1 2 3 4 5
\end{verbatim}

\paragraph{standard output}
\begin{verbatim}
10
55
\end{verbatim}

\end{document}
